\begin{graphillustration}{Пример: свёртка КСС даёт конденсацию. Нумерация справа — топологический порядок; все дуги идут слева направо.}{fig:03-condensation}
\begin{tikzpicture}[font=\small]
\node[gi vertex] (a) at (0,0) {$a$};
\node[gi vertex] (b) at (1.1,0) {$b$};
\node[gi vertex] (c) at (3,0) {$c$};
\node[gi vertex] (d) at (4.1,0) {$d$};
\node[gi vertex] (e) at (6,0) {$e$};
\draw[gi arrow] (a) to[bend left=25] (b);
\draw[gi arrow] (b) to[bend left=25] (a);
\draw[gi arrow] (c) to[bend left=25] (d);
\draw[gi arrow] (d) to[bend left=25] (c);
\draw[gi arrow] (b) -- (c);
\draw[gi arrow] (d) -- (e);
\draw[gi arrow] (a) to[bend left=30] (e);
\node[gi group,fit=(a)(b),label=below:$C_1$] {};
\node[gi group,fit=(c)(d),label=below:$C_2$] {};
\node[gi group,fit=(e),label=below:$C_3$] {};
\node[gi note] at (8.2,0) {$\longrightarrow$};
\begin{scope}[shift={(9.3,0)}]
\node[gi vertex] (C1) at (0,0) {$C_1$};
\node[gi vertex] (C2) at (1.4,0) {$C_2$};
\node[gi vertex] (C3) at (2.8,0) {$C_3$};
\draw[gi arrow] (C1) -- (C2);
\draw[gi arrow] (C2) -- (C3);
\draw[gi arrow] (C1) to[bend left=35] (C3);
\end{scope}
\node[gi note] at (10.7,-1.3) {$\varphi(C_1)=1,\ \varphi(C_2)=2,\ \varphi(C_3)=3$};
\end{tikzpicture}
\end{graphillustration}
